\scanpage{108}

Pour définir un iso

\[
X^\Omega\xrightarrow{\sim}\hat\Sigma
\]

il revient au $=$, par descente, de définir un iso

\[
(X^\Omega)_\Omega\xrightarrow{\sim}(\hat\Sigma)_\Omega
\]

compatible avec les $\Omega$-autom des deux schém., ou ce qui revient au $=$ [comme

\[
(X^\Omega)_\Omega\simeq(X_\Omega)^{\Omega_\Omega=\Omega\times_S\Omega=\Omega\amalg\Omega}\simeq X_\Omega\times X_\Omega,
\]

avec l'autom de descente (provenant de l'indice $\Omega$ à $X^\Omega$) se transformant en \inserted{le} l'autom$_{\!S}$ de $X_\Omega\times X_\Omega$ composé de la symétrie des 2 facteurs et de l'autom can. au dessus de l'autom non trivial $\sigma_\Omega$ de $\Omega$, il revient au $=$ de trouver un isom.

\[
X_\Omega\times X_\Omega\xrightarrow{\sim}\hat\Sigma_\Omega\quad(=\widehat{\Sigma_\Omega})
\]

qui transforme $\sigma_\Omega^{\unsure{(X\times X)_\Omega}}\timess$ en $\sigma_\Omega^{(\hat\Sigma_\Omega)}$.

Cela signifie aussi : pour tt $\omega\in\Gamma(\unsure{S'},(\det V)^{\unsure{\otimes2}})$ ($S'$ quelc sur $S$) tel que

\[
\delta'(q\deleted{\unsure{_{S'}}})=\omega^{\otimes-2}
\]

trouver un iso

\[
f_\omega:X\times X\longrightarrow\hat\Sigma\tag*{$(*)$}
\]

\deleted{\unsure{can.}} satisfaisant les conditions

% archive sheet number: 107
